% ============================================================
%  Chapter 1 -- Introduction
% ============================================================
\chapter{Introduction}
\label{ch:intro}

\section{Motivation}

Three-dimensional geometric data has grown rapidly in the past decade, driven
by advances in acquisition, storage, and transmission.
LiDAR sensors on autonomous vehicles stream dense point clouds at tens of
frames per second.
Volumetric video capture systems record human bodies as time-varying 3D
surfaces for immersive telepresence and extended reality applications.
Cultural heritage institutions digitise artefacts at millimetre precision
for permanent archival.
In all of these settings, the point cloud --- an unordered set of 3D
coordinates, optionally augmented with colour or normal attributes ---
has become the standard representation: it is directly acquired, imposes no
topological constraints, and avoids the cubic memory cost of dense voxel grids.

Efficient compression is a prerequisite for any practical use of this data.
Raw point clouds are large: a single frame of a human body capture at MPEG
standard resolution contains up to 800,000 points, each requiring three
floating-point coordinates.
At video frame rates, uncompressed storage or transmission is prohibitive.
The MPEG standards body has responded with two families of codecs ---
G-PCC~\cite{graziosi2020gpcc} for static or sparsely sampled geometry, and
V-PCC~\cite{schwarz2018emerging} which projects 3D surfaces onto 2D patches
and applies conventional video compression --- but neither achieves the
rate-distortion efficiency that learned methods now routinely demonstrate on
image and video benchmarks.

Learned point cloud codecs, led by architectures such as
PCGCv2~\cite{wang2021multiscale} and its successors, apply the
analysis-synthesis paradigm from neural image compression directly to 3D
occupancy grids.
An encoder maps the point cloud to a compact latent representation; a learned
entropy model estimates the latent distribution for arithmetic coding;
a decoder inverts the process.
At high bitrates these systems match or exceed G-PCC.
At low bitrates, however, they reveal a consistent and under-studied failure
mode: \emph{geometric oversmoothing}.

\begin{figure}[t]
  \centering
  \includegraphics[width=0.92\textwidth]{soldier_vis.png}
  \caption{%
    Geometric oversmoothing in learned point cloud compression (PCGCv2 at r2,
    0.05~bpp).
    \textbf{Left}: original \textit{soldier} sequence, frame 803 (79,529 pts).
    \textbf{Centre}: PCGCv2 decoded (76,058 pts) --- geometric detail at surface
    boundaries and protrusions is smoothed away by the learned encoder.
    \textbf{Right}: proposed post-processed output (76,058 pts) --- fine structure
    substantially recovered.
    The oversmoothing artifact is invisible to standard D1 PSNR, which averages
    uniformly over all surface points.
  }
  \label{fig:teaser}
\end{figure}

\section{The Oversmoothing Problem}

Oversmoothing in learned codecs arises from a fundamental property of entropy
models: rare, fine-grained occupancy patterns incur high coding cost.
High-curvature regions --- edges, ridges, creases, corners --- produce exactly
such patterns in the occupancy grid.
At low bitrates, the encoder is implicitly penalised for representing these
regions faithfully, and the decoder reconstructs a surface in which such
features are blurred or erased.

The consequence is spatially non-uniform degradation: flat surfaces are
reconstructed near-perfectly while geometric detail concentrates the error.
Standard quality metrics such as point-to-point PSNR (D1)~\cite{tian2017geometric}
average over the entire reconstructed surface, weighting every point equally
regardless of its geometric role.
A codec that faithfully reconstructs 90\% of a surface while systematically
erasing all edges can still report a high aggregate PSNR --- the metric masks
precisely the information that matters perceptually.

Figure~\ref{fig:teaser} illustrates this effect: the compressed cloud looks
reasonable at a glance but has lost all fine surface geometry at boundaries
and protrusions.
This thesis begins from the observation that measuring and correcting
oversmoothing requires evaluating quality \emph{where it hurts}: at
high-curvature regions of the original surface.

\section{Scope and Contributions}

This work makes the following contributions:

\begin{enumerate}[label=\textbf{\arabic*.}]
  \item \textbf{Curvature-stratified PSNR metric.}
    We introduce an evaluation framework that partitions the original point
    cloud by local surface curvature and computes PSNR separately for edge
    and flat strata.
    The metric uses the reverse (original-to-reconstructed) nearest-neighbour
    direction, which measures how well the reconstruction covers the original
    surface rather than the converse.
    Curvature is estimated from the eigenvalue structure of the local
    covariance matrix; stratification is adaptive via percentile thresholds,
    making it dataset-independent.

  \item \textbf{Systematic characterisation of oversmoothing.}
    We apply the metric to PCGCv2 across four sequences of the
    8i~Voxelized Full Bodies (8iVFB) dataset at seven operating rates.
    Edge-region PSNR consistently lags flat-region PSNR by 2--4.5~dB, and
    the gap \emph{widens} with increasing rate --- the codec allocates
    additional bits toward flat surfaces, not toward recovering detail.

  \item \textbf{Rate-conditioned sparse U-Net post-processor.}
    We propose a lightweight post-processing network that refines decoded point
    positions without modifying the codec.
    The backbone is a three-level sparse U-Net operating directly on the
    decoded occupancy grid.
    A Feature-wise Linear Modulation (FiLM)~\cite{perez2018film} head adapts
    displacement predictions to the compressed bitrate, enabling a single model
    to serve multiple operating points.
    Training uses a dynamic Chamfer Distance loss that circumvents the
    zero-inflation and nearest-neighbour noise problems that defeat
    displacement-regression losses.

  \item \textbf{Experimental validation and generalisation.}
    The model achieves $+2.06$~dB edge PSNR at the lowest evaluated rate
    (r2), $+0.58$~dB at medium rate (r4), and $+0.27$~dB at high rate (r6),
    averaged across all four 8iVFB sequences and 400 frames.
    Zero-shot transfer to an unseen codec (Unicorn) and an unseen dataset
    (MVUB) demonstrates that low-rate oversmoothing is a cross-codec failure
    mode and that the correction generalises without retraining.
\end{enumerate}

\section{Thesis Organisation}

The remainder of this thesis is structured as follows.

\textbf{Chapter~\ref{ch:background}} develops the theoretical foundations
necessary to understand the work: the evolution from 1D signal processing to
3D geometric processing, the properties of point clouds and their compression,
the deep learning tools used (sparse convolutions, autoencoders, FiLM
conditioning), and the quality metrics employed throughout.

\textbf{Chapter~\ref{ch:oversmoothing}} presents the curvature-stratified
evaluation framework, applies it to PCGCv2 and G-PCC, and quantifies the
oversmoothing pattern that motivates the post-processing approach.

\textbf{Chapter~\ref{ch:model}} describes the post-processing model in detail:
problem formulation, backbone architecture, FiLM conditioning, loss function
design, and training procedure.

\textbf{Chapter~\ref{ch:results}} reports experimental results including ablation
studies on the loss function, curvature features, and rate weighting strategies,
full 400-frame evaluation on 8iVFB, and zero-shot generalisation experiments.

\textbf{Chapter~\ref{ch:conclusion}} summarises the findings, discusses
limitations, and outlines directions for future work.
